% © 2026 D. Jaroš — CC BY-NC-ND 4.0
%
% File: research_tracks/T3_ALPHA/alpha_weinberg_final_status.tex
% Purpose: Final status summary of alpha and Weinberg angle derivations.
%          All proof levels, all remaining conditions, all bridge assumptions.
%
% Date: 2026-06-14

% UBT-AI-PROVENANCE-BEGIN
% schema: ubt-ai-provenance/v1
% tier: C_working
% ai_assistance: disclosed
% human_review: risk-based
% editorial_responsibility: Ing. David Jaroš
% policy: ../../AI_PROVENANCE.md
% notice: Working material; exhaustive human review is not claimed.
% UBT-AI-PROVENANCE-END

\ifdefined\INCLUDEMODE\else
\documentclass[12pt,a4paper]{article}
\usepackage[a4paper,margin=2.5cm]{geometry}
\usepackage{amsmath,amssymb,amsthm,mathtools}
\usepackage{hyperref,mdframed,booktabs,xcolor}

\newtheorem{theorem}{Theorem}[section]
\newtheorem{lemma}[theorem]{Lemma}
\newtheorem{proposition}[theorem]{Proposition}
\newtheorem{corollary}[theorem]{Corollary}
\theoremstyle{remark}\newtheorem{remark}[theorem]{Remark}

\newmdenv[backgroundcolor=green!15,linecolor=green!60!black,linewidth=1.5pt]{resultbox}
\newmdenv[backgroundcolor=yellow!20,linecolor=orange,linewidth=1.5pt]{gapbox}
\newmdenv[backgroundcolor=blue!8,linecolor=blue!50,linewidth=1.5pt]{insightbox}

\newcommand{\Lz}{\textbf{[L0]}}
\newcommand{\Lo}{\textbf{[L1]}}
\newcommand{\Lc}{\textbf{[L1 cond.]}}
\newcommand{\MC}{\textbf{[MC]}}
\newcommand{\STD}{\textbf{[STD]}}

\title{\textbf{Final Status: $\alpha^{-1}$ and $\sin^2\theta_W$}\\[4pt]
\large Proof levels, bridge assumptions, and remaining conditions}
\author{D.~Jaros}
\date{2026-06-14}

\begin{document}
\maketitle
\fi

%──────────────────────────────────────────────────────────────────
\section{Fine-Structure Constant $\alpha^{-1} = 137.035999177549$}
\label{sec:alpha}
%──────────────────────────────────────────────────────────────────

\begin{resultbox}
\textbf{Result.}
$\alpha^{-1}_{\rm UBT} = 137.035999177549\ldots$,
$+0.026\sigma$ from CODATA/NIST 2022. No fitted parameter.

\textbf{Overall status}: \Lc\ (all bridge assumptions stated below).
\end{resultbox}

\subsection{Mathematical chain (all \Lo\ or \Lz)}

\begin{center}
\renewcommand{\arraystretch}{1.25}
\begin{tabular}{lll}
\toprule
Result & Level & Source \\
\midrule
$\Omega_\eta(1)=\tfrac{1}{24}-\tfrac{1}{8\pi}$ & \Lz & Eisenstein $G_2(i)$ \\
$12\pi\Omega_\eta = (\pi-3)/2$ & \Lz & algebra \\
$N_{\rm eff}=12=3\times4$ & \Lo & \texttt{su2\_twist\_neff12.tex} \\
$Z_\psi[\tau]=\eta^{-12}$ (exact Gaussian) & \Lo & \texttt{modular\_symmetry\_rpsi.tex} \\
$R_\psi=1$ (modular fixed point) & \Lo & \texttt{modular\_symmetry\_rpsi.tex} \\
$r_{\rm L2}=3(1+\Omega_\eta)$ & \Lo & \texttt{layer2\_kernel\_derivation.tex} \\
$\rho=\exp(-C_Q/(2\pi n))$ (exact Gaussian) & \Lo & \texttt{gap\_A\_proof.tex} \\
$n^*=137.035999177549$ (fixed point) & \Lo & \texttt{alpha\_derivation\_complete.tex} \\
\bottomrule
\end{tabular}
\end{center}

\subsection{Physical identification $n^* = \alpha^{-1}$ (\Lc)}

The UBT EM action on the $\psi$-winding sector is
\[
  S^{\rm UBT}_{\rm EM}[A]
  = \frac{n}{8\pi}\int F\wedge *F - i\frac{\chi}{4\pi}\int F\wedge F,
\]
giving complexified coupling $\tau^{\rm UBT}_{\rm EM} = \chi + in$.
Matching to the standard Maxwell coupling
$\tau_{\rm EM} = \theta_{\rm EM}/(2\pi) + i\cdot4\pi/e^2$:
\[
  \frac{n}{8\pi} = \frac{1}{2e^2}
  \;\;\Longrightarrow\;\;
  n = \frac{4\pi}{e^2} = \alpha^{-1}. \quad \Lc
\]
[\texttt{em\_modular\_coupling\_identification.tex}]

\subsection{Remaining conditions for unconditional \Lo}

\begin{gapbox}
\begin{enumerate}
\item \textbf{EM projection condition} \Lc:
  The UBT compact-$\psi$ sector must project onto the standard
  unit-charge $U(1)_{\rm EM}$ line bundle.
  Requires: $\Theta_0$ vacuum orbit in the EW doublet with $Y=1$
  [\texttt{electroweak\_em\_projection\_theorem.tex}].

\item \textbf{EW non-zero minimum} \Lc:
  $V(\Theta)$ must have $\mu^2_{\rm EW}>0$ (SSB condition).
  Status: proved conditionally on the UBT-compatible potential form
  [\texttt{theta0\_nonzero\_minimum\_theorem.tex}].
  Open: derive $\mu^2_{\rm EW}>0$ from $S[\Theta]$ dynamics.
\end{enumerate}
\end{gapbox}

%──────────────────────────────────────────────────────────────────
\section{Weinberg Angle $\sin^2\theta_W(M_Z) = 0.23122$}
\label{sec:weinberg}
%──────────────────────────────────────────────────────────────────

\begin{resultbox}
\textbf{Result.}
$\sin^2\theta_W(M_Z) = 0.23122\ldots$,
matching experiment $0.23122(3)$ to $5$ significant figures.

\textbf{Overall status}: \Lc\ within stated bridge assumptions.
\end{resultbox}

\subsection{Derivation chain}

\textbf{Boundary condition} [\Lo]:
From hypercharge normalisation in UBT:
\[
  \sin^2\theta_W(M_{\rm UBT}) = \frac{3}{8}.
\]

\textbf{Beta coefficients} [\Lc, \texttt{weinberg\_angle\_odd\_spinor\_threshold\_beta\_theorem.tex}]:
From odd-spinor threshold spectrum (paired EW doublets, Grassmann measure,
3-generation structure):
\[
  b_1 = \frac{33}{5}, \quad b_2 = 1, \quad b_3 = -3.
\]

\textbf{RG running} [\STD]:
\[
  \alpha_i^{-1}(M_Z) = A + b_i L,
  \qquad
  L = \frac{1}{2\pi}\ln\frac{M_{\rm UBT}}{M_Z} = 5.2555\ldots
\]
Hypercharge: $\alpha_Y^{-1} = \tfrac{5}{3}\alpha_1^{-1}$.
EM assembly:
\[
  \alpha_{\rm EM}^{-1}(M_Z)
  = \tfrac{8}{3}A + 12L = 127.934\ldots \quad [\STD\ {\rm input}]
\]
\[
  \Longrightarrow\quad
  A = \frac{3}{8}\bigl[\alpha_{\rm EM}^{-1}(M_Z) - 12L\bigr]
  = 24.325\ldots
\]
\[
  \sin^2\theta_W(M_Z)
  = \frac{\alpha_2^{-1}(M_Z)}{\alpha_{\rm EM}^{-1}(M_Z)}
  = \frac{A + L}{127.934}
  = \frac{29.581}{127.934}
  = 0.23122\ldots \quad \checkmark
\]

\subsection{Bridge assumptions and sensitivity}

\begin{gapbox}
\textbf{Bridge B5} \MC: Physical scale $M_{\rm UBT} \approx 2\times10^{16}$~GeV.
This is the primary open bridge. Sensitivity: $\sin^2\theta_W$ shifts by
$\pm0.003$ per order of magnitude in $M_{\rm UBT}$.
Derivation of $M_{\rm UBT}$ from $S[\Theta]$: open.

\textbf{Bridge B9} \Lc: $\alpha_{\rm EM}^{-1}(M_Z) = 127.934$ from
$\alpha_{\rm UBT}^{-1}=137.036$ plus standard QED running.
\STD\ input used directly.

\textbf{Odd-spinor threshold} \Lc: beta coefficients require
paired EW doublets and Grassmann measure from odd-winding sector.
[\texttt{weinberg\_angle\_odd\_spinor\_threshold\_beta\_theorem.tex}]
\end{gapbox}

%──────────────────────────────────────────────────────────────────
\section{Combined Status Matrix}
\label{sec:matrix}
%──────────────────────────────────────────────────────────────────

\begin{center}
\renewcommand{\arraystretch}{1.3}
\begin{tabular}{lllp{4.5cm}}
\toprule
Quantity & Value & Level & Key open condition \\
\midrule
$\alpha^{-1}$ & $137.035999177549$ & \Lc & EW doublet vacuum ($\Theta_0$) \\
$\sin^2\theta_W(M_Z)$ & $0.23122$ & \Lc\ (B5) & Physical scale $M_{\rm UBT}$ \\
$M_{\rm UBT}$ & $\approx 2\times10^{16}$~GeV & \MC & Derive from $S[\Theta]$ \\
$\mu^2_{\rm EW}>0$ & (SSB sign) & \Lc & Derive from $S[\Theta]$ dynamics \\
\bottomrule
\end{tabular}
\end{center}

\begin{insightbox}
\textbf{What is and is not claimed.}

The mathematical fixed-point chain giving $n^*=137.036$ is \Lo\ (exact,
no approximation). The physical identification $n^*=\alpha^{-1}$ via
compact-U(1) EM normalisation is \Lc. The Weinberg angle derivation is
\Lc\ inside stated bridge assumptions.

Neither result should be advertised as unconditionally proved.
The correct statement is: \emph{``within the UBT bridge framework,
$\alpha^{-1}=137.036$ and $\sin^2\theta_W(M_Z)=0.231$ follow from
first principles with all derivation steps at \Lo\ or \Lc\ level.''}

The two primary remaining tasks to promote to unconditional \Lo:
\begin{enumerate}
\item Derive $M_{\rm UBT}$ (physical scale) from $S[\Theta]$ (B5).
\item Derive $\mu^2_{\rm EW}>0$ (SSB sign) from $S[\Theta]$ dynamics.
\end{enumerate}
\end{insightbox}

\begin{thebibliography}{9}
\bibitem{alpha_complete}
  D.~Jaros, \texttt{alpha\_derivation\_complete.tex}, 2026.
\bibitem{em_id}
  D.~Jaros, \texttt{em\_modular\_coupling\_identification.tex}, 2026.
\bibitem{ew_proj}
  D.~Jaros, \texttt{electroweak\_em\_projection\_theorem.tex}, 2026.
\bibitem{theta0_min}
  D.~Jaros, \texttt{theta0\_nonzero\_minimum\_theorem.tex}, 2026.
\bibitem{odd_beta}
  D.~Jaros, \texttt{weinberg\_angle\_odd\_spinor\_threshold\_beta\_theorem.tex}, 2026.
\bibitem{weinberg_closure}
  D.~Jaros, \texttt{weinberg\_angle\_mz\_bridge\_closure.tex}, 2026.
\bibitem{bridge_matrix}
  D.~Jaros, \texttt{BRIDGE\_CLOSURE\_MATRIX.md}, 2026.
\end{thebibliography}

\ifdefined\INCLUDEMODE\else
\end{document}
\fi
